\begin{definition}[The energy-normalized gradient operator $D$]
For $v \in W^{1,2}(\gamma_n)$ in the sense of \noderef{SobolevW12}, let
\[
  Dv = 2^{-1/2}\nabla v \in L^2(\gamma_n;\mathbb{R}^n),
\]
that is, the unique element $\Gamma$ of $L^2(\gamma_n;\mathbb{R}^n)$ such that
$\sqrt2\,\Gamma$ is a weak Euclidean gradient of $v$ in the sense of
\noderef{HasGaussWeakGrad}; uniqueness as an element of
$L^2(\gamma_n;\mathbb{R}^n)$ is \noderef{GaussWeakGradUnique}.  Its $j$th
coordinate is $D_j v = 2^{-1/2}\partial_j v$, with $\partial_j$ understood in the
distributional sense.  For $v \in L^2(\gamma_n)\setminus W^{1,2}(\gamma_n)$ we set
$Dv = 0$, a convention that makes $D$ a total map on $L^2(\gamma_n)$ and that is
never invoked in any statement below: every statement mentioning $Dv$ carries the
hypothesis that $v$ lies in $W^{1,2}(\gamma_n)$ or in a subspace of it.
\end{definition}
